\documentclass[11pt]{article}
\usepackage[margin=0.9in]{geometry}
\usepackage{enumitem}
\usepackage{amssymb}
\usepackage{xcolor}
\usepackage[T1]{fontenc}
\usepackage{titlesec}

\definecolor{accent}{HTML}{4B2E83}
\titleformat{\section}{\large\bfseries\color{accent}}{}{0em}{}[\titlerule]
\titlespacing*{\section}{0pt}{10pt}{5pt}

\newcommand{\cmd}[1]{\texttt{\small #1}}
\newlist{mycheck}{itemize}{1}
\setlist[mycheck]{label=$\square$, leftmargin=1.6em, itemsep=2pt, topsep=3pt}

\pagestyle{empty}

\begin{document}

{\LARGE\bfseries\color{accent} Windows Defense Checklist}\\[2pt]
{\large CCDC Tryout Prep \textbullet{} Windows Client and Server}\\[2pt]
\textit{Rule of thumb: before you disable or change anything, ask ``could a scored service depend on this?'' Back it up first, and write down every change and why.}

\section*{0. First Five Minutes}
\begin{mycheck}
  \item Read the scenario and any README or rules fully; note which services must stay up
  \item Open a notes file and log each change with a timestamp (tryouts also evaluate communication)
  \item Identify the OS version, hostname, and role (client, member server, domain controller, IIS, file server, database)
  \item Take a snapshot or backup before major changes, if allowed
  \item \textbf{Back up services before touching them} (full commands in the appendix). Create \cmd{C:\textbackslash backup} first, outside any share and outside \cmd{inetpub}
  \begin{itemize}[itemsep=1pt, topsep=1pt]
    \item \textbf{Shares:} run \cmd{Get-SmbShare}, note each Name and Path (ignore hidden defaults like C\$, ADMIN\$, IPC\$), zip each shared folder, and export share names and permissions to CSV
    \item \textbf{IIS:} copy \cmd{C:\textbackslash inetpub} (no space in the path) and run \cmd{Backup-WebConfiguration}
    \item \textbf{MySQL:} find \cmd{mysqldump.exe} (often not in PATH) and dump with \cmd{--result-file}
  \end{itemize}
  \item Practice a restore at least once in your lab; a backup you have never restored is a guess
\end{mycheck}

\section*{1. Accounts and Passwords}
\begin{mycheck}
  \item Change default and weak passwords (Administrator and any service accounts); for unauthorized users, just change the password
  \item Ask your team how to handle scoring service accounts before changing them
  \item Review accounts and groups in \cmd{lusrmgr.msc}; only needed service accounts and authorized users should exist
  \item List local users: \cmd{Get-LocalUser}
  \item List Administrators group: \cmd{Get-LocalGroupMember Administrators}
  \item Accounts should only be in groups they are authorized for
  \item On a DC, review Domain Admins, Enterprise Admins, Schema Admins: \cmd{Get-ADGroupMember "Domain Admins"}
  \item Disable Guest and the default Administrator account if not needed
  \item Make sure passwords expire; check for accounts with ``never expires'' or ``no password required''
  \item Set password and lockout policy (\cmd{secpol.msc}, \cmd{gpmc.msc}, or a PowerShell script)
  \item Remember domain logins use \cmd{DOMAIN\textbackslash user}: the part before the backslash is what you are signing into
\end{mycheck}

\section*{2. Processes, Services, and Network}
\begin{mycheck}
  \item Review running processes: \cmd{Get-Process}, \textbf{Process Explorer} (check parent/child chains and injected DLLs)
  \item Use \textbf{Process Monitor} for real-time file, registry, and process activity
  \item List services and startup types: \cmd{Get-Service}; manage with \cmd{sc} (create, delete, stop, start, configure)
  \item Check \textbf{Server Manager} (\cmd{servermanager.exe}) for service status and best-practice warnings
  \item Identify which services are scored before stopping anything
  \item Stop and disable clearly unneeded services (e.g., unused remote access tools)
  \item List listening ports: \cmd{Get-NetTCPConnection -State Listen}; map each port to a process
  \item Run \textbf{TCPView} (Sysinternals) to see live connections per process; watch for odd outbound connections
  \item Use \textbf{Everything} (run as admin) and search \cmd{*.exe} to find unexpected executables
\end{mycheck}

\section*{3. Persistence}
\begin{mycheck}
  \item Run \textbf{Autoruns} (Sysinternals); review Logon, Services, Scheduled Tasks, and Drivers tabs. Red entries are suspicious, but Autoruns does not show everything
  \item Check \textbf{Task Scheduler} / \cmd{Get-ScheduledTask} for attacker-created tasks
  \item Check for \textbf{IFEO (image hijacks)}: look for \cmd{Debugger} values under the Image File Execution Options registry key (e.g., sethc.exe replaced with cmd.exe); delete them
  \item Check \textbf{PowerShell profiles} (\cmd{profile.ps1}); they run every time PowerShell opens. If a command mysteriously fails, an alias may be set to null or to something malicious
  \item Check startup folders and Run/RunOnce registry keys (HKLM and HKCU)
  \item Look for unexpected WMI subscriptions and new services
  \item Check the \textbf{root certificate store} for planted certificates; if you find a malicious cert, see whether other malware uses it
  \item Remember attackers abuse signed Windows services or sign their own binaries, so ``signed'' does not mean safe
  \item Check \cmd{appwiz.cpl} for unwanted software; in File Explorer enable View $>$ Show $>$ Hidden items and File name extensions
  \item Look for unknown files in user profiles, Temp, and ProgramData
  \item Forensics: \textbf{ShellBags Viewer} shows folder view settings stored in the registry
\end{mycheck}

\section*{3a. Using Autoruns (Sysinternals)}
\textit{Autoruns lists nearly everything that runs at startup or logon. It does not show everything, so use it with Task Scheduler, services, and Process Explorer.}
\begin{mycheck}
  \item \textbf{Get it:} download from the Microsoft Sysinternals site, or use a tools folder your team provides. Ask your coach whether downloads and internet are allowed. Use \cmd{autoruns64.exe} (GUI) and \cmd{autorunsc64.exe} (command line, good for Server Core)
  \item \textbf{Run as administrator} (right-click, Run as administrator), or some entries will be missing. Accept the license prompt
  \item \textbf{Save a baseline first:} File $>$ Save, early on a clean-looking machine. Later use File $>$ Compare to see only \textbf{new} entries (shown in green)
  \item Start on the \textbf{Everything} tab, then check these tabs: \textbf{Logon}, \textbf{Services}, \textbf{Scheduled Tasks}, \textbf{Drivers}, \textbf{Image Hijacks}, \textbf{WMI}, \textbf{Winlogon}, \textbf{AppInit}, \textbf{Boot Execute}
  \item Turn on Options $>$ Scan Options $>$ \textbf{Verify code signatures}. \textbf{Check VirusTotal.com} sends file hashes to the internet, so use it only if you have internet and it is allowed
  \item \textbf{Colors:} pink means no publisher or an invalid signature; yellow means the entry exists but its file is missing; green means new since your saved baseline. With VirusTotal enabled, red detection counts are bad. Your team's ``red is bad'' refers to these suspicious highlights, so investigate every highlighted row
  \item \textbf{Hide Microsoft Entries} reduces noise, but attackers abuse signed Windows programs. Also look at the full list for \cmd{cmd.exe}, \cmd{powershell.exe}, \cmd{mshta.exe}, \cmd{rundll32.exe}, or \cmd{wscript.exe} with odd arguments, or entries pointing to Temp, AppData, ProgramData, Public, or Users folders
  \item \textbf{Suspicious signs:} no publisher, random file names, script files (.ps1, .vbs, .bat), encoded PowerShell commands, files in odd folders, drivers you do not recognize
  \item Click an entry to see its path, publisher, and command line in the bottom pane. Right-click for \textbf{Jump to Entry} (registry location), \textbf{Jump to Image} (file location), \textbf{Properties}, and \textbf{Search Online}
  \item \textbf{Disable before deleting:} uncheck an entry to disable it (reversible), then re-test your scored services. Delete it only once you are confident it is malicious. Never remove entries for scored services (IIS, MySQL, DNS, and so on)
  \item Delete \textbf{Image Hijacks} (IFEO) entries pointing to \cmd{cmd.exe} or other programs
  \item If you find a malicious signed file, check the certificate and see whether other files use it. Cross-check with Task Scheduler, \cmd{Get-Service}, and Process Explorer
  \item Re-run Autoruns (F5 to refresh) regularly; attackers re-add persistence
  \item Log what you found: entry name, location, path, and what you did
\end{mycheck}
{\footnotesize
\begin{verbatim}
cmd /c "autorunsc64.exe -accepteula -a * -c -h -s -m -nobanner > C:\backup\autoruns.csv"
\end{verbatim}
}
\noindent The command above saves a CSV of non-Microsoft entries with file hashes and signature checks. Run it again later and compare the files.

\section*{4. Firewall and Shares}
\begin{mycheck}
  \item Open \cmd{wf.msc} (Windows Defender Firewall with Advanced Security; built on the Windows Filtering Platform). Confirm it is on for all profiles: \cmd{Get-NetFirewallProfile}
  \item Typical approach: delete default rules, add only what is needed, set default inbound and outbound to block. \textbf{Do this only after} you have listed what the scored services require (Section 7)
  \item Remember: connections that do not match an allow rule are blocked, so apps like Firefox need an outbound program rule (New Rule $\rightarrow$ Program $\rightarrow$ path to the .exe)
  \item Allow only required inbound ports (e.g., 80 for IIS, 445 for SMB)
  \item Use predefined rules (e.g., RDP) and tighten them in Properties, such as limiting remote IPs; or use \cmd{New-NetFirewallRule}
  \item \textbf{simplewall} (free, third party) is \emph{not} a front end for Windows Firewall; it configures the Windows Filtering Platform directly and runs independently alongside it. If either firewall blocks a connection, it is blocked, so check \textbf{both} when a service breaks. Keep Windows Firewall on as a second layer
  \item simplewall is whitelist-style: it blocks all apps by default, so you allow only what is needed. Its filters stay active even after simplewall is closed; the \textbf{Disable filters} button removes them
  \item \textbf{simplewall workflow:}
  \begin{itemize}[itemsep=1pt, topsep=1pt]
    \item Finish backups and list each scored service's ports and programs first (Section 7 baseline)
    \item Open simplewall but \textbf{do not enable filters yet}; create your allow rules first
    \item Prefer \textbf{temporary} rules while learning (reset on reboot), or \cmd{simplewall.exe -install -temp}
    \item Enable filters, then test every service from another machine after each rule
    \item Know your rollback before you start: Disable filters, or reboot if rules are temporary
  \end{itemize}
  \item Custom rules are \textbf{global} (all apps) or \textbf{special} (chosen apps, set on the rule's Apps tab). Services run by the System process (IIS, SMB) are likely easier to allow with a \textbf{port-based} rule than a per-program rule; verify this in your lab
  \item \textbf{Inbound allow, by machine role:} IIS TCP 80 (443 if HTTPS is scored); SMB TCP 445; DNS TCP and UDP 53; MySQL TCP 3306 (limit to the web server IP if possible); RDP TCP 3389 only if needed, ideally from specific IPs; domain controller needs many AD ports (53, 88, 135, 389, 445, 464, 636, 3268/3269, dynamic RPC), so get the exact list from your team first
  \item \textbf{Outbound allow, only what the machine starts itself:} DNS (53) to your DNS server; DHCP if used; NTP (123) for AD time; web server to MySQL on 3306; domain members to the DC; Windows Update only if allowed and safe. Replies to allowed inbound connections need no extra rule
  \item Ask your team for their exact simplewall setup and copy it; practice in your lab until you can enable, test, and disable it quickly
  \item A \textbf{share} is a folder other machines reach over SMB (port 445), e.g., \cmd{\textbackslash\textbackslash SERVER\textbackslash Docs}; the scoring engine may test it. List shares with \cmd{Get-SmbShare} or \cmd{fsmgmt.msc} (check Name and Path; ignore hidden defaults C\$, ADMIN\$, IPC\$)
  \item Remove unneeded shares and check share and NTFS permissions; \textbf{never share the C: drive}
  \item Keep your backup folder (\cmd{C:\textbackslash backup}) \textbf{unshared} and outside \cmd{inetpub}
  \item If a share is deleted, recreate it with \cmd{New-SmbShare -Name Docs -Path C:\textbackslash Shares\textbackslash Docs} and restore permissions from your exported CSV
  \item Restrict or disable RDP unless required; require NLA if kept
  \item Disable SMBv1 if not needed
\end{mycheck}

\section*{5. Active Directory, DNS, and Group Policy (Servers)}
\begin{mycheck}
  \item Verify AD, DNS, and DHCP services are healthy and responding
  \item Check \cmd{dnsmgmt.msc} for unexpected records; DNS is a very common place for competition breaks
  \item Review GPOs in \cmd{gpmc.msc} (domain) or \cmd{gpedit.msc} (single computer) for suspicious or recently changed settings
  \item Check for unexpected delegated permissions and new OUs or GPO links
  \item Confirm domain-joined machines can still authenticate after changes
  \item Know how to manage \textbf{Server Core} (no GUI) with PowerShell and \cmd{sconfig}
\end{mycheck}

\section*{6. Web and Application Services (IIS, MySQL)}
\begin{mycheck}
  \item Open IIS Manager (\cmd{inetmgr}); review sites, bindings, application pools, and server certificates/SSL settings
  \item IIS breaks easily: check \textbf{bindings}, \textbf{handler mappings}, \textbf{read permissions} on site files, and the inbound firewall rule for 0.0.0.0:80
  \item Look for unexpected files in web roots (\cmd{C:\textbackslash inetpub}), including web shells
  \item Remove default or sample content; disable directory browsing
  \item MySQL: back up with \cmd{mysqldump}; check users and remote access
  \item When something breaks, check in order: the service, firewall, permissions, bindings/config, then DNS
\end{mycheck}

\section*{7. Keeping Services Up (First 2 Hours)}
\textit{Goal: keep every scored service running and fix what breaks. Change low-risk things first (passwords, firewall cleanup); hold off on risky hardening until you know what is scored.}
\begin{mycheck}
  \item Get the list of scored services and which machine hosts each one
  \item Record a baseline: listening ports per service (\cmd{Get-NetTCPConnection -State Listen}) and the firewall rules each one needs
  \item Test every service from \textbf{another machine}, the way the scoring engine sees it: \cmd{Test-NetConnection HOST -Port 80}
  \item \textbf{Troubleshooting order, every time:} (1) service running? (2) listening on the right port? (3) firewall allows it? (4) file/account permissions? (5) config and bindings? (6) DNS resolving?
  \item \textbf{IIS:} site and app pool started; bindings (IP, port, hostname); handler mappings; read permissions on the site folder; authentication and default document; inbound rule for port 80
  \item \textbf{SMB:} Server service running; share still exists (\cmd{Get-SmbShare}); share and NTFS permissions; port 445 open
  \item \textbf{MySQL:} service running and listening on 3306; firewall; user accounts and grants; restore from your dump if data was dropped
  \item \textbf{DNS:} zones and A records correct in \cmd{dnsmgmt.msc}; forwarders; test with \cmd{Resolve-DnsName} or \cmd{nslookup}
  \item \textbf{Active Directory:} run \cmd{dcdiag}; DNS healthy (AD depends on it); system time correct; NTDS and Netlogon services running
  \item Find what changed before blindly fixing: Event Viewer, Task Scheduler, GPOs, recently modified files
  \item Make \textbf{one change at a time} and re-test after each
  \item Restore the service first, then investigate. If something malicious is causing the break, fix that cause too
  \item Log every break: what broke, the cause, the fix (injects often want a written report)
\end{mycheck}

\section*{8. Updates, Logging, and Cleanup}
\begin{mycheck}
  \item Ensure Windows Defender Antivirus is on and definitions are current (separate from the firewall)
  \item Install critical updates when allowed and safe for scored services
  \item Enable auditing for logons, account changes, and process creation
  \item Review Event Viewer (\cmd{eventvwr.msc}): Security (\textbf{4624 successful login}, \textbf{4625 failed login}, 4720, 4732), System, Application, PowerShell logs
  \item A flood of 4625 events suggests brute force; unexpected 4624 events suggest a compromised account
  \item Remove unneeded software and tools
  \item Advanced: enable \textbf{PPL} (Protected Process Light) for LSASS via registry so DLLs cannot be injected
  \item Be careful in the registry (\cmd{regedit}); one wrong key can break a lot. Advanced system settings: \cmd{sysdm.cpl}
\end{mycheck}

\section*{9. Verify and Document}
\begin{mycheck}
  \item Re-test every scored service after your changes
  \item Re-check for new accounts, tasks, and connections (attackers return, especially in the second half facing the red team)
  \item Finish your change log: what, why, when
  \item Be ready to explain your reasoning and what you would do next
\end{mycheck}

\newpage
{\LARGE\bfseries\color{accent} Appendix: PowerShell Quick Reference}\\[2pt]
\textit{Replace names, paths, and IPs with your own. Run PowerShell as Administrator. Test commands in your lab before using them in competition.}

\section*{Accounts}
{\footnotesize
\begin{verbatim}
Get-LocalUser | Select Name,Enabled,PasswordExpires,PasswordRequired
Set-LocalUser -Name bob -Password (Read-Host -AsSecureString)
Disable-LocalUser -Name Guest
Get-LocalGroupMember Administrators
Remove-LocalGroupMember -Group Administrators -Member bob
Get-ADUser -Filter * -Properties PasswordNeverExpires | Where PasswordNeverExpires
Set-ADAccountPassword -Identity bob -Reset -NewPassword (Read-Host -AsSecureString)
Get-ADGroupMember "Domain Admins"
\end{verbatim}
}

\section*{Services and Processes}
{\footnotesize
\begin{verbatim}
Get-Service W3SVC,LanmanServer,MySQL,DNS,NTDS | ft Name,Status,StartType
Start-Service W3SVC;  Restart-Service W3SVC
Set-Service -Name SomeSvc -StartupType Disabled
sc.exe qc SomeSvc                       # show a service's config and binary path
Get-Process | Sort CPU -Descending | Select -First 15
Get-CimInstance Win32_Process | Select Name,ProcessId,ParentProcessId,CommandLine
Get-AuthenticodeSignature C:\path\file.exe
Get-WinEvent -FilterHashtable @{LogName='System';Level=2;StartTime=(Get-Date).AddHours(-1)}
\end{verbatim}
}

\section*{Network and Firewall}
{\footnotesize
\begin{verbatim}
Get-NetTCPConnection -State Listen | Select LocalAddress,LocalPort,OwningProcess
Get-NetTCPConnection -State Listen | Select LocalPort,
    @{n='Proc';e={(Get-Process -Id $_.OwningProcess).Name}}
Get-NetTCPConnection -State Established
Test-NetConnection 10.0.0.5 -Port 80
Get-NetFirewallProfile | Select Name,Enabled,DefaultInboundAction,DefaultOutboundAction
Get-NetFirewallRule -Direction Inbound -Enabled True | Select DisplayName,Action
New-NetFirewallRule -DisplayName "Allow HTTP" -Direction Inbound -Protocol TCP `
    -LocalPort 80 -Action Allow
Disable-NetFirewallRule -DisplayName "Rule Name"
Set-NetFirewallProfile -Profile Domain,Private,Public -DefaultInboundAction Block
\end{verbatim}
}

\section*{Shares and SMB}
{\footnotesize
\begin{verbatim}
New-Item -ItemType Directory C:\backup -Force     # not shared, not inside inetpub
Get-SmbShare | Select Name,Path,Special           # ignore Special (C$, ADMIN$, IPC$)
Get-SmbShareAccess -Name ShareName

# Zip one share (use the Path from Get-SmbShare)
Compress-Archive -Path C:\Shares\Docs -DestinationPath C:\backup\Docs.zip

# Zip every non-default share, and save share names and permissions
Get-SmbShare | Where-Object { -not $_.Special } | ForEach-Object {
    Compress-Archive -Path $_.Path -DestinationPath "C:\backup\$($_.Name).zip" -Force
}
Get-SmbShare | Export-Csv C:\backup\shares.csv -NoTypeInformation
Get-SmbShare | Get-SmbShareAccess | Export-Csv C:\backup\share-access.csv -NoTypeInformation

# If zip fails (large folder or files in use), copy and keep NTFS permissions
robocopy C:\Shares\Docs C:\backup\Docs /E /COPYALL

# Restore
Expand-Archive C:\backup\Docs.zip -DestinationPath C:\Shares
New-SmbShare -Name Docs -Path C:\Shares\Docs       # recreate a deleted share

Set-SmbServerConfiguration -EnableSMB1Protocol $false -Force
\end{verbatim}
}

\section*{IIS}
{\footnotesize
\begin{verbatim}
Import-Module WebAdministration
Get-Website;  Get-WebBinding
Get-ChildItem IIS:\AppPools
Start-Website "Default Web Site";  Start-WebAppPool DefaultAppPool
icacls C:\inetpub\wwwroot              # check read permissions
Backup-WebConfiguration -Name pre-change          # sites, bindings, app pools
Copy-Item C:\inetpub C:\backup\inetpub -Recurse   # no space in the path
robocopy C:\inetpub C:\backup\inetpub /E /COPYALL # alternative; keeps permissions
iisreset
\end{verbatim}
}

\section*{MySQL}
{\footnotesize
\begin{verbatim}
Get-Command mysqldump -ErrorAction SilentlyContinue    # is it in PATH?
# If not, check C:\Program Files\MySQL\MySQL Server 8.0\bin or C:\xampp\mysql\bin
$bin = "C:\Program Files\MySQL\MySQL Server 8.0\bin"    # adjust to your version
& "$bin\mysqldump.exe" -u root -p --all-databases --result-file=C:\backup\all.sql
Get-Item C:\backup\all.sql                             # confirm it exists, not empty

# Restore (PowerShell has no <, so go through cmd)
cmd /c '"C:\Program Files\MySQL\MySQL Server 8.0\bin\mysql.exe" -u root -p < C:\backup\all.sql'
mysql -u root -p -e "SELECT user,host FROM mysql.user;"
\end{verbatim}
}
\noindent Use \cmd{--result-file} instead of \cmd{>} in PowerShell, which can write the dump in the wrong encoding. \cmd{-p} prompts for the password, so it stays out of your command history.

\section*{DNS and Active Directory}
{\footnotesize
\begin{verbatim}
Get-DnsServerZone
Get-DnsServerResourceRecord -ZoneName corp.local
Resolve-DnsName www.corp.local -Server 127.0.0.1
Add-DnsServerResourceRecordA -Name www -ZoneName corp.local -IPv4Address 10.0.0.5
Export-DnsServerZone -Name corp.local -FileName corp.local.bak
dcdiag /q
repadmin /replsummary
w32tm /query /status
Get-GPO -All | Select DisplayName,ModificationTime
\end{verbatim}
}

\section*{Persistence Hunting}
{\footnotesize
\begin{verbatim}
Get-ScheduledTask | Where State -ne Disabled | Select TaskName,TaskPath
Get-CimInstance Win32_StartupCommand | Select Name,Command,Location
Get-ItemProperty HKLM:\Software\Microsoft\Windows\CurrentVersion\Run
Get-ItemProperty HKCU:\Software\Microsoft\Windows\CurrentVersion\Run
$k = "HKLM:\SOFTWARE\Microsoft\Windows NT\CurrentVersion"
Get-ItemProperty "$k\Image File Execution Options\*" | Where Debugger   # IFEO hijacks
Get-CimInstance -Namespace root/subscription -ClassName __EventConsumer
Get-Content $PROFILE.AllUsersAllHosts   # also check $PROFILE for each user
Get-ChildItem Cert:\LocalMachine\Root | Sort NotBefore -Descending |
    Select -First 10 Subject,Thumbprint,NotBefore
Get-ChildItem C:\Users -Recurse -Include *.exe -Force -ErrorAction SilentlyContinue
\end{verbatim}
}

\section*{Event Logs}
{\footnotesize
\begin{verbatim}
Get-WinEvent -FilterHashtable @{LogName='Security';Id=4625} -MaxEvents 20   # failed logins
Get-WinEvent -FilterHashtable @{LogName='Security';Id=4624} -MaxEvents 20   # successful logins
\end{verbatim}
}

\section*{Advanced}
{\footnotesize
\begin{verbatim}
# Enable PPL for LSASS (reboot required; test in your lab first)
Set-ItemProperty HKLM:\SYSTEM\CurrentControlSet\Control\Lsa -Name RunAsPPL -Value 1 -Type DWord
sconfig                                 # Server Core configuration menu
\end{verbatim}
}

\vfill
\small\textit{Tip: practice in a lab (Windows Server VM with IIS, SMB, DNS, MySQL; break one thing at a time, then fix it) until you can run this routine from memory, and prefer PowerShell over clicking through the GUI. Time yourself on Sections 0--4 and 7.}

\end{document}
